\subsection*{\textbf{Cuestionario}}

\subsubsection*{\textit{1. ¿Cuáles son los principales beneficios de la utilización de procedimientos almacenados en la gestión de bases de datos?}}

Los procedimientos almacenados ofrecen múltiples beneficios fundamentales en la administración de bases de datos. \parencite{Base100} destaca que la \textbf{optimización del rendimiento} es uno de los principales, ya que la ejecución precompilada de estos procedimientos acelera significativamente las operaciones. Además, se logra un \textbf{ahorro de espacio} al evitar la repetición de código mediante la centralización de la lógica en un único lugar. \parencite{Bosco2020} resalta que el \textbf{tráfico de red reducido} entre cliente y servidor es otra ventaja crítica, puesto que los comandos de un procedimiento se ejecutan en un único lote de código. La \textbf{reutilización del código} permite que múltiples aplicaciones accedan a la misma lógica sin necesidad de reescribir código redundante. \parencite{Gonzales2022} enfatiza el \textbf{mantenimiento más sencillo}, ya que cuando se necesitan modificaciones en los procesos de base de datos, solo debe actualizarse la lógica en la capa de datos, dejando independiente el nivel de aplicación. Finalmente, la \textbf{reducción de errores} se minimiza al usar código probado y centralizado.

\subsubsection*{\textit{2. ¿Cuáles son los distintos tipos de procedimientos almacenados?}}

Existen varios tipos de procedimientos almacenados clasificados según su alcance y ubicación. \parencite{Base100} y \parencite{Bosco2020} describen los siguientes tipos: Los \textbf{procedimientos almacenados del sistema} se encuentran guardados en la base de datos maestra y permiten a los administradores recuperar información de las tablas del sistema. Los \textbf{procedimientos almacenados locales} son creados individualmente en las bases de datos de cada usuario. Los \textbf{procedimientos almacenados temporales} pueden ser locales (disponibles solo durante una sesión de usuario) o globales (disponibles para todas las sesiones). Los \textbf{procedimientos almacenados remotos} solo son manejados por algunos DBMS como SQL Server, permitiendo la ejecución en servidores distintos. Los \textbf{procedimientos almacenados extendidos} se implementan como librerías DLL, se ejecutan fuera del entorno del DBMS pero se invocan como procedimientos normales. \parencite{MySQLDocumentation} complementa esta clasificación indicando que en MySQL los procedimientos también pueden clasificarse según su visibilidad y uso específico en aplicaciones.

\subsubsection*{\textit{3. ¿Cuál es la definición técnica de un Trigger SQL disparador en el contexto de una base de datos relacional?}}

Un trigger o disparador es \parencite{Bosco2020} ``una especie de procedimiento almacenado que se ejecuta automáticamente cuando ocurre un evento sobre alguna tabla''. \parencite{Walther2023} define más específicamente que un trigger es un objeto de base de datos que se activa automáticamente en respuesta a eventos específicos (INSERT, UPDATE, DELETE) en tablas particulares. A diferencia de los procedimientos almacenados, los triggers no pueden recibir parámetros de entrada del usuario ni devolver valores explícitamente; en su lugar, operan sobre los datos que están siendo modificados. \parencite{MySQLDocumentation} indica que los triggers se definen con un timing específico (BEFORE o AFTER) y un evento específico, permitiendo ejecutar lógica compleja de validación, auditoría o automatización en respuesta a cambios en los datos. Los triggers actúan como guardianes de la integridad de datos y automatizadores de procesos sin requerer intervención explícita del usuario.

\subsubsection*{\textit{4. ¿En qué situaciones o casos de uso se recomienda la implementación de un trigger?}}

Los triggers son especialmente útiles en múltiples escenarios de aplicación en bases de datos. \parencite{Walther2023} describe la \textbf{automatización de tareas} como un caso de uso primario; por ejemplo, cada vez que se registra un nuevo usuario, un trigger puede actualizar automáticamente una tabla de auditoría. La \textbf{integridad de datos} es otro caso crítico donde los triggers pueden aplicar reglas que protejan la consistencia; por ejemplo, impidiendo eliminar registros que aún tienen referencias en otras tablas. \parencite{Bosco2020} enfatiza el uso de triggers para \textbf{auditoría y registro histórico}, mantener un registro de cambios en datos específicos incluyendo quién cambió qué y cuándo. Las \textbf{acciones en cascada} representan otro caso de uso importante; si se elimina un registro principal, un trigger puede eliminar automáticamente registros relacionados en tablas secundarias. \parencite{Base100} complementa indicando que los triggers son ideales para \textbf{validación compleja de datos}, ejecutando lógica de negocios que va más allá de las restricciones simples de base de datos, permitiendo reglas condicionales sofisticadas.

\subsubsection*{\textit{5. ¿Qué evento o condicción provoca la activación automática de un trigger?}}

Los triggers se activan automáticamente en respuesta a tres eventos fundamentales en SQL. \parencite{Bosco2020} especifica que estos eventos son: \textbf{INSERT} (inserción de registros nuevos), \textbf{UPDATE} (actualización de registros existentes) y \textbf{DELETE} (eliminación de registros). \parencite{Celaya} complementa que cada una de estas acciones puede definir desencadenadores dentro de la tabla especificada. El timing de la activación es crítico: \parencite{MySQLDocumentation} indica que los triggers pueden configurarse como \textbf{BEFORE} (ejecutándose antes de que el evento ocurra, permitiendo prevención o modificación) o \textbf{AFTER} (ejecutándose después del evento, permitiendo auditoría o acciones secundarias). Según \parencite{Walther2023}, la especificidad del trigger permite que solo ciertos tipos de cambios activen la lógica del disparador; por ejemplo, un trigger puede estar configurado para activarse solo ante inserciones específicas que cumplan ciertos criterios dentro del bloque de código.

\subsubsection*{\textit{6. ¿Cuál es la sintaxis general o el comando utilizado para crear un trigger SQL?}}

La sintaxis general para crear un trigger en SQL sigue una estructura específica. \parencite{Bosco2020} proporciona la sintaxis fundamental:

\begin{verbatim}
DELIMITER //
CREATE TRIGGER nombre_trigger
BEFORE | AFTER INSERT | UPDATE | DELETE
ON nombre_tabla
FOR EACH ROW
BEGIN
    sentencias SQL
END //
DELIMITER ;
\end{verbatim}

\parencite{Celaya} aclara que es necesario modificar el delimitador (DELIMITER) de su valor predeterminado semicolon (;) antes de crear el trigger. El delimitador personalizado permite al sistema distinguir entre el final de cada sentencia SQL dentro del trigger y el final de la definición completa del trigger. \parencite{MySQLDocumentation} completa la sintaxis indicando que después de BEFORE o AFTER va el evento (INSERT, UPDATE, DELETE), seguido de ON con el nombre de la tabla. La cláusula \texttt{FOR EACH ROW} especifica que el trigger se ejecute para cada fila afectada. Finalmente, el bloque BEGIN...END contiene las sentencias SQL que se ejecutarán.

\subsubsection*{\textit{7. ¿Cuáles son los distintos tipos de triggers disparadores que existen, generalmente clasificados por el momento de ejecución o el tipo de evento?}}

Los triggers pueden clasificarse de múltiples formas según su timing y evento asociado. \parencite{Bosco2020} clasifica por momento de ejecución: los \textbf{BEFORE triggers} se ejecutan antes de que ocurra la acción, permitiendo validar datos, prevenir inserciones/actualizaciones inválidas o modificar los datos antes de guardarlos. Los \textbf{AFTER triggers} se ejecutan después de que la acción ha ocurrido, siendo ideales para auditoría, registros de cambios, o efectuar acciones secundarias dependientes de cambios confirmados. \parencite{Walther2023} clasifica por tipo de evento: los \textbf{INSERT triggers} se activan cuando se insertan nuevos registros, \textbf{UPDATE triggers} cuando se modifican registros existentes, y \textbf{DELETE triggers} cuando se eliminan registros. Es importante notar que \parencite{MySQLDocumentation} indica que se pueden definir múltiples triggers para la misma tabla y evento, aunque se ejecutarán en el orden en que fueron creados. La combinación de timing y evento resulta en seis tipos principales: BEFORE INSERT, AFTER INSERT, BEFORE UPDATE, AFTER UPDATE, BEFORE DELETE, y AFTER DELETE.

\subsubsection*{\textit{8. ¿Considerando su funcionamiento, admiten los triggers la recepcción de parámetros de entrada o salida, a diferencia de los procedimientos almacenados?}}

No, los triggers \textbf{no admiten parámetros de entrada o salida} explícitos, a diferencia de los procedimientos almacenados. \parencite{Bosco2020} clarifica que ``un TRIGGER no puede devolver datos al usuario, sino que simplemente se ejecuta para cambiar o comprobar los datos que se van a insertar, actualizar o borrar''. \parencite{Celaya} complementa señalando que aunque los triggers operan automáticamente en respuesta a eventos, no aceptan argumentos del usuario como lo hacen los procedimientos. Sin embargo, \parencite{MySQLDocumentation} destaca que los triggers tienen acceso automático a referencias especiales: las variables \textbf{NEW} y \textbf{OLD}. \parencite{Bosco2020} explica que \textbf{NEW} referencia los nuevos valores de los campos después de la modificación (disponible en INSERT y UPDATE), mientras que \textbf{OLD} referencia los valores anteriores (disponible en UPDATE y DELETE). De esta forma, aunque no aceptan parámetros convencionales, los triggers tienen acceso contextual a los datos siendo modificados, permitiendo lógica condicional basada en estos valores implícitos.

\subsubsection*{\textit{9. ¿Cuáles son los requisitos esenciales o las consideraciones fundamentales que deben tenerse en cuenta al diseñar e implementar un trigger?}}

Existen varios requisitos y consideraciones críticas al implementar triggers según \parencite{Bosco2020}: Primero, \textbf{las tablas asociadas deben existir} antes de crear el trigger. Segundo, \textbf{el nombre del trigger no debe ser una palabra reservada} del sistema de base de datos. Tercero, es necesario \textbf{modificar el delimitador de ejecución} de su valor predeterminado (;) a uno personalizado usando DELIMITER, aunque técnicamente no es obligatorio, es una práctica recomendada. \parencite{Base100} advierte que \textbf{un trigger solo se puede aplicar a una tabla}, aunque puede referenciar objetos fuera de su base de datos. Aunque un trigger se defina dentro de una sola base de datos, puede hacer referencia a objetos que se encuentren fuera de la misma. \parencite{MySQLDocumentation} precisa que \textbf{no se permiten sentencias DDL} (CREATE, ALTER, DROP) dentro de la definición de un trigger. Además, \parencite{Bosco2020} señala que \textbf{los triggers no pueden utilizarse en vistas}, solo en tablas base. Finalmente, es importante considerar el \textbf{rendimiento y la recursión}; los triggers pueden afectar el rendimiento si contienen lógica compleja, y algunos sistemas limitan la recursión de triggers para prevenir bucles infinitos.

\subsubsection*{\textit{10. ¿Los triggers SQL están limitados a su uso en sistemas de gestión de bases de datos específicos como Oracle o SQL Server o son una característica estándar de SQL?}}

Los triggers son una característica ampliamente soportada pero con variaciones entre diferentes DBMS. \parencite{SQLStandard} indica que los triggers están definidos en el estándar SQL:2016, convirtiéndose en una característica estándar de SQL que la mayoría de sistemas modernos implementan. Sin embargo, \parencite{Bosco2020} y \parencite{Base100} subrayan que existen \textbf{diferencias de sintaxis y capacidades} entre sistemas. \parencite{MySQLDocumentation} documenta que MySQL implementa triggers con soporte para BEFORE/AFTER INSERT/UPDATE/DELETE, pero con limitaciones en comparación con Oracle o SQL Server. Por ejemplo, \parencite{Bosco2020} menciona que algunos sistemas como SQL Server soportan procedimientos almacenados remotos que otros DBMS no implementan. \parencite{Walther2023} confirma que prácticamente todos los sistemas relacionales modernos (MySQL, PostgreSQL, Oracle, SQL Server) soportan triggers, aunque sus sintaxis específicas y características avanzadas pueden variar. La conclusión es que mientras los triggers son estándares en SQL y disponibles en la mayoría de DBMS, su implementación exacta, limitaciones y características avanzadas varían según el sistema específico, requiriendo que los desarrolladores comprendan las particularidades de su sistema objetivo.
